% §11.1 Miller-Rabin Formal
%
% Recasts the chapter 3 §3.4 informal walkthrough as a formal theorem about
% a $\Bernoulli{\BoolType}$ instance. Retires the §3.4 deferral by stating
% the per-round and $k$-round error bounds as theorems with proper citation
% to the number-theoretic underpinnings (Rabin 1980).
%
% Cross-reference resolution status:
%
% RESOLVED (chapters 3 and 4, drafted in Plans 1-3):
%   \Cref{sec:miller-rabin}         -- chapter 3 §3.4 informal walkthrough
%   \Cref{thm:composition}          -- chapter 3 §3.3 composition theorem
%   \Cref{sec:axiom-1}              -- chapter 4 §4.2 element independence
%   \Cref{ax:element-indep}         -- Axiom 1 formal statement
%   \Cref{def:bernoulli-bool}       -- chapter 4 §4.5 $\Bernoulli{\BoolType}$
%
% INTERNAL (resolves later in chapter 11):
%   \Cref{sec:sample-complexity}    -- chapter 11 §11.4 (Plan 11 Task 5)

\section{Miller-Rabin Formal}
\label{sec:miller-rabin-formal}

Chapter 3 \S3.4 walked through the Miller-Rabin primality test as an
extended example, showing step by step how a single round behaves as a
$\Bernoulli{\BoolType}$ instance with one-sided error and how $k$ rounds
compose by the chain rule of \Cref{thm:composition}. The exposition there
was deliberately informal in one specific respect: the per-round bound
$\fprate \le 1/4$ was stated as a fact about strong-liar sets and then
applied. The number-theoretic origin of that bound was acknowledged but
deferred. This section closes the deferral. We restate the §3.4 walkthrough
(see \Cref{sec:miller-rabin}) as a single formal theorem with an explicit
appeal to Rabin's original result, then derive the $k$-round bound as a
corollary of \Cref{thm:composition}.

%---------------------------------------------------------------------------
\subsection*{The formal statement}

Fix an odd integer $n > 3$. Write $n - 1 = 2^r d$ with $d$ odd and $r \ge 1$.
For each candidate witness $a \in \{2, 3, \ldots, n-2\}$, let
$\mathrm{MillerRabinRound}(n, a)$ denote the single-round witness test of
\Cref{sec:miller-rabin}: it returns \textsc{composite} when $a$ certifies
the compositeness of $n$ via the square-roots-of-1 criterion, and
\textsc{probably-prime} otherwise. Define the witness predicate
\begin{equation*}
  W_n(a) \;:=\; \mathrm{MillerRabinRound}(n, a),
\end{equation*}
viewed as a Boolean-valued function on the witness pool
$\{2, \ldots, n-2\}$ (with \textsc{probably-prime} identified with $\top$
and \textsc{composite} with $\bot$, as in chapter 3 \S3.4 Step 2).

\begin{theorem}[Miller-Rabin as a $\Bernoulli{\BoolType}$]
\label{thm:miller-rabin-formal}
Let $n > 3$ be an odd integer, and let $a$ be drawn uniformly at random
from $\{2, \ldots, n-2\}$. The single-round witness predicate $W_n(a)$ is
a $\Bernoulli{\BoolType}$ instance (in the sense of
\Cref{def:bernoulli-bool}) whose error rates satisfy
\begin{equation}
\label{eq:mr-formal-fnr}
  \fnrate \;=\; 0
\end{equation}
for every $n$, and
\begin{equation}
\label{eq:mr-formal-fpr}
  \fprate \;\le\; \frac{1}{4}
\end{equation}
for every odd composite $n > 3$.

If $k$ independent witnesses $a_1, \ldots, a_k$ are drawn uniformly with
replacement from $\{2, \ldots, n-2\}$, and the $k$-round test declares
\textsc{probably-prime} only when every $W_n(a_i)$ does, then the $k$-round
error rates satisfy
\begin{equation}
\label{eq:mr-formal-k-round}
  \fnrate^{(k)} \;=\; 0
  \quad\text{and}\quad
  \fprate^{(k)} \;\le\; \left(\frac{1}{4}\right)^{\!k}
\end{equation}
for every odd composite $n > 3$.
\end{theorem}

%---------------------------------------------------------------------------
\subsection*{Proof sketch}

The two claims of \Cref{thm:miller-rabin-formal} have different status. The
first claim, $\fprate \le 1/4$ per round, is a number-theoretic fact about
the structure of $\mathbb{Z}/n\mathbb{Z}$. The second claim, the $k$-round
bound, is a structural consequence that the Bernoulli framework supplies
once the first is in hand. We sketch each in turn.

\begin{proof}[Proof sketch of \Cref{thm:miller-rabin-formal}]
\emph{Zero false-negative rate.} Suppose $n$ is prime. For any
$a \in \{2, \ldots, n-2\}$, Fermat's little theorem gives
$a^{n-1} \equiv 1 \pmod{n}$. The sequence
$a^d, a^{2d}, a^{4d}, \ldots, a^{2^r d} \pmod{n}$ terminates at $1$, and
in $\mathbb{Z}/n\mathbb{Z}$ (which is a field when $n$ is prime) the only
square roots of $1$ are $\pm 1$. Tracing the sequence backwards from
$a^{2^r d} = 1$, the last value before $1$ must itself be $\pm 1$. The
witness test passes. Hence $W_n(a) = \textsc{probably-prime}$ for every
witness whenever $n$ is prime, which gives $\fnrate = 0$.

\smallskip
\emph{Per-round false-positive bound.} Suppose $n$ is odd composite. Define
the \emph{strong-liar set} for $n$,
\[
  L_n \;=\; \{\, a \in \{2, \ldots, n-2\} : W_n(a) = \textsc{probably-prime} \,\}.
\]
Rabin's theorem~\cite{rabin-1980}, building on the earlier deterministic
analysis of Miller~\cite{miller-1976}, establishes the strong-liar bound
\[
  |L_n| \;\le\; \frac{n-3}{4}.
\]
This is the number-theoretic content of Miller-Rabin; the proof goes
through the structure of $\mathbb{Z}/n\mathbb{Z}$ via the Chinese Remainder
Theorem and a careful counting of solutions to $x^{2^j d} \equiv \pm 1
\pmod{n}$ for composite $n$. We do not reproduce it here; \S11.6 records
the standard secondary references. With $|L_n| \le (n-3)/4$
and $a$ drawn uniformly from a pool of size $n - 3$,
\[
  \fprate \;=\; \Pr[W_n(a) = \textsc{probably-prime} \mid n \text{ composite}]
  \;=\; \frac{|L_n|}{n - 3} \;\le\; \frac{1}{4},
\]
which is the bound of equation~\eqref{eq:mr-formal-fpr}.

\smallskip
\emph{The $k$-round bound.} The argument here is exactly the chapter 3
\S3.4 Step 3 calculation, now phrased as an application of the framework's
Axiom 1 (\Cref{ax:element-indep}, \Cref{sec:axiom-1}) followed by the
composition theorem. Recast the per-round test as a predicate
$\tilde{P}\colon \{2, \ldots, n-2\} \to \{\top, \bot\}$ on the witness
pool, so that distinct witnesses $a_i \neq a_j$ are distinct input
elements. Drawing the $a_i$ independently and uniformly makes the events
$\{a_i \in L_n\}$ mutually independent (this is sampling-with-replacement
independence, which holds by the protocol, not by any deeper property of
$L_n$). Axiom 1 therefore applies to $\tilde{P}$. The $k$-round test is
the $k$-fold AND of $\tilde{P}(a_1), \ldots, \tilde{P}(a_k)$, an
identity-with-error chain in the sense of \Cref{thm:composition} with
$\fnrate_i = 0$ and $\fprate_i \le 1/4$ for each $i$.

Substituting into the composition theorem gives
\[
  \fnrate^{(k)}
  \;=\; 1 - \prod_{i=1}^{k}(1 - \fnrate_i)
  \;=\; 1 - 1
  \;=\; 0,
\]
and, in the FNR-zero special case where errors at later rounds cannot
cancel an earlier false positive, the false-positive rate of the chain
reduces to the product
\[
  \fprate^{(k)}
  \;\le\; \prod_{i=1}^{k} \fprate_i
  \;\le\; \left(\frac{1}{4}\right)^{\!k},
\]
which is equation~\eqref{eq:mr-formal-k-round}. The bound is uniform over
all odd composite $n > 3$, since the per-round bound was.
\qedhere
\end{proof}

The proof has the shape that chapter 3 \S3.4 promised. One ingredient is
borrowed from number theory and credited to its origin
\cite{miller-1976,rabin-1980}; the rest is the Bernoulli framework applied
mechanically. The choice of $k$ to drive $\fprate^{(k)}$ below a target
failure probability $\delta$ is a sample-complexity question that we take
up in \Cref{sec:sample-complexity}.

%---------------------------------------------------------------------------
\subsection*{What changed from \S3.4}

The §3.4 walkthrough and the present section are about the same algorithm,
the same per-round bound, and the same $k$-round bound. What changed is the
status of the per-round bound in the surrounding argument.

In §3.4, the bound $\fprate \le 1/4$ was introduced as a fact to be
applied. We named it the Miller-Rabin theorem and attached it to the
strong-liar set, but we did not prove it and we did not cite the original
source. The pedagogical reason was clear: §3.4 was Part~I climax 1, whose
purpose was to demonstrate that the chapter 3 composition machinery could
carry a real algorithm. Pausing to prove a 1980 number-theory theorem would
have distracted from the framework story. The bound was a black box, and
the framework was what filled the rest of the page.

In §11.1, the same bound is stated as part of \Cref{thm:miller-rabin-formal}
and credited to its origins~\cite{rabin-1980,miller-1976}. The proof sketch
makes the division of labor explicit: number theory supplies the
strong-liar bound, the framework supplies everything else. We have not
strengthened the §3.4 result, nor have we recomputed it. What we have done
is reframe the §3.4 walkthrough as the proof sketch of a single formal
theorem, with the per-round bound cited rather than asserted, and the
$k$-round bound derived as a corollary of \Cref{thm:composition} rather
than as a self-contained product calculation.

This reframing is the contribution of \S11.1, and it is small by design.
The chapter 11 promise is not to re-derive Miller-Rabin from scratch; it
is to retire the §3.4 informal layer by replacing it with a formal layer
that names its borrowed pieces. The remaining sections of chapter 11
extend the same kind of formalization to the broader landscape of
randomized algorithms: the Las Vegas versus Monte Carlo distinction, the
complexity classes RP and BPP, and the sample-complexity bounds that
govern how many repetitions are needed to reach a target error rate.
